\chapter{Impacto del trabajo} \label{chp:impacto}
Este capítulo está dedicado a evaluar las repercusiones y los beneficios derivados del desarrollo de este Trabajo de Fin de Grado. En la primera sección se abordará el impacto general del proyecto desde una perspectiva tecnológica, organizativa y socioeconómica, destacando cómo las decisiones de diseño tomadas han estado condicionadas por la necesidad de aportar un valor real y seguro al entorno laboral. En la segunda sección, se analizará la alineación de la plataforma desarrollada con los Objetivos de Desarrollo Sostenible (ODS) de la Agenda 2030 de las Naciones Unidas.
%%%%%%%%%%%%%%%%%%%%%%%%%%%%%%%%%%%%%%%%%%%%%%%%%%%%%%%%%%%%%%%%%%%%%%%%%%%%%%%% %%%%%%%%%%%%%%%%%%%%%%%%%%%%%%%%%%%%%%%%%%%%%%%%%%%%%%%%%%%%%%%%%%%%%%%%%%%%%%%%
\section{Impacto general} \label{sct:impacto:general}
El desarrollo de este sistema automatizado de altas de empleados presenta un impacto potencial altamente positivo, especialmente en el contexto organizativo de las gestorías y los departamentos de Recursos Humanos. Hasta ahora, estos sectores se han visto obligados a lidiar con procesos burocráticos repetitivos y dependientes de software de escritorio anticuado (como la aplicación SILTRA) debido a las limitaciones tecnológicas del Sistema RED de la Seguridad Social.
La implementación de esta plataforma aporta una modernización drástica, transformando un proceso propenso a errores humanos en un flujo digital, ágil y guiado. Al incorporar validaciones algorítmicas en tiempo real (como la verificación de identificadores DNI/NIE o el control lógico de fechas de nacimiento y alta), se reduce significativamente el riesgo de enviar datos erróneos a la administración, lo cual a su vez previene posibles sanciones administrativas para las empresas y ahorra horas de trabajo de corrección de expedientes.
A lo largo del desarrollo, se han tomado diversas decisiones técnicas fundamentadas en maximizar este impacto positivo y mitigar riesgos:
\begin{itemize} \item \textbf{Adopción de una arquitectura Multi-Tenant:} La decisión más crítica fue el diseño de la base de datos documental en Firestore orientada a múltiples inquilinos. Esto se hizo con la consideración innegociable de proteger el derecho a la privacidad. Al gestionar datos altamente sensibles (identificación, salarios, domicilios), la arquitectura garantiza una segregación estricta de la información, de modo que ningún usuario pueda acceder a datos de empresas para las que no está explícitamente autorizado, cumpliendo así de forma rigurosa con el RGPD. \item \textbf{Uso de tecnologías Low-Code y Backend-as-a-Service:} La elección de FlutterFlow y Firebase se basó en el impacto de la escalabilidad y el mantenimiento. Se democratiza el desarrollo de software corporativo permitiendo a las gestorías, que normalmente no disponen de grandes presupuestos informáticos, acceder a una infraestructura en la nube robusta, segura y fácil de mantener. \item \textbf{Desarrollo de una Mock API de transición:} Ante la imposibilidad técnica de conectarse a una API pública moderna de la Tesorería General de la Seguridad Social (TGSS), se decidió crear un sistema de validación local propio. El impacto de esta decisión es proporcionar un entorno de pruebas y una arquitectura ya preparada (Plug and Play) para que, en un futuro, la plataforma pueda engancharse directamente a una pasarela comercial u oficial de forma transparente para el usuario final. \end{itemize}
%%%%%%%%%%%%%%%%%%%%%%%%%%%%%%%%%%%%%%%%%%%%%%%%%%%%%%%%%%%%%%%%%%%%%%%%%%%%%%%% %%%%%%%%%%%%%%%%%%%%%%%%%%%%%%%%%%%%%%%%%%%%%%%%%%%%%%%%%%%%%%%%%%%%%%%%%%%%%%%%
\section{Objetivos de Desarrollo Sostenible} \label{sct:impacto:ods}
El proyecto desarrollado no solo busca una mejora tecnológica, sino que también contribuye de forma directa y transversal a la consecución de varios Objetivos de Desarrollo Sostenible (ODS) impulsados por la Agenda 2030 de la ONU. Las metas más relevantes que se ven impactadas por la naturaleza de este trabajo son las siguientes:
\begin{itemize}
\item \textbf{ODS 8 - Trabajo decente y crecimiento económico} \ Este objetivo promueve el crecimiento económico sostenido, inclusivo y sostenible, el empleo pleno y productivo, y el trabajo decente. La aplicación desarrollada tiene un impacto directo en este ODS al facilitar y agilizar la formalización legal de los contratos laborales. Al simplificar la burocracia para dar de alta a los trabajadores en la Seguridad Social, se combate la economía sumergida y se fomenta la contratación regulada, garantizando que los empleados accedan a sus derechos sociales desde el primer día de actividad. Además, aumenta la productividad y la eficiencia económica de las pequeñas y medianas empresas (pymes) y despachos laborales, permitiéndoles enfocar sus recursos en tareas de mayor valor en lugar de en la burocracia repetitiva.
\item \textbf{ODS 9 - Industria, innovación e infraestructura} \ La meta central de este ODS es construir infraestructuras resilientes, promover la industrialización inclusiva y sostenible y fomentar la innovación. El proyecto aporta una innovación tecnológica clave al sector de la administración laboral, que tradicionalmente ha sufrido un rezago tecnológico. Al introducir paradigmas de desarrollo como el \textit{Low-Code}, las bases de datos NoSQL e integraciones en la nube (Cloud Functions), la plataforma dota a las gestorías de una infraestructura digital moderna, escalable y accesible que moderniza su tejido industrial y sus procesos diarios.
\item \textbf{ODS 16 - Paz, justicia e instituciones sólidas} \ Este ODS aboga por crear instituciones eficaces, responsables y transparentes a todos los niveles. Al haber diseñado la aplicación con un enfoque férreo en el cumplimiento de la legalidad (integrando validaciones requeridas por el SEPE y la TGSS) y en la seguridad de la información (garantizando el Reglamento General de Protección de Datos - RGPD mediante una separación de roles y arquitectura \textit{Multi-Tenant}), el sistema fortalece la relación de transparencia y confianza entre las empresas, los trabajadores y la Administración Pública. Asimismo, el registro histórico (colección \textit{altas\_history}) garantiza una trazabilidad y auditoría completa de los trámites, previniendo irregularidades y fomentando la responsabilidad institucional.
\end{itemize}

